% Omslagmotief wiskunde: assenstelsel met een sinusoïde, een e-macht met raaklijn en een cirkel die de sinus in de
% top raakt. Alles is exact: de punten liggen echt op de grafieken (nagerekend, zie de getallen hieronder).
% Tekenvlak ca. x -8..118, y -8..58 (mm); onder y -8 begint de titel. De oorsprong ligt op (0,12).
\begin{scope}[shift={(0,12)}]
  % assen
  \draw[white,opacity=.35,line width=.6pt,->] (-8,0) -- (118,0);
  \draw[white,opacity=.35,line width=.6pt,->] (0,-18) -- (0,46);
  % y = 12 sin(6x), periode 60: toppen bij x = 15 en 75, dalen bij x = 45 en 105 (y = -12, net boven de titel)
  \draw[kleur1!60!white,line width=1.6pt,domain=-8:116,samples=200] plot (\x,{12*sin(6*\x)});
  % y = 2,2 e^(x/20)
  \draw[kleur2!70!white,line width=1.6pt,domain=-8:58,samples=120] plot (\x,{2.2*exp(\x/20)});
  % raaklijn in x = 40: y = 2,2e^2 = 16,256, helling y/20 = 0,8128
  \draw[white,opacity=.55,line width=.7pt] (28,6.502) -- (52,26.010);
  \fill[white] (40,16.256) circle (.9);
  % snijpunten van sinus en e-macht (numeriek): x = 1,943 en x = 23,817
  \fill[white] (1.943,2.425) circle (.75);
  \fill[white] (23.817,7.238) circle (.75);
  % cirkel met straal 15 die de sinus in de top (75, 12) raakt; straal 15 > kromtestraal 7,6, dus alleen daar
  \draw[kleur3!55!white,line width=1.6pt] (75,27) circle (15);
  \fill[white] (75,27) circle (.9);
  \draw[white,opacity=.6,line width=.6pt,dashed] (75,27) -- (85.607,37.607);
  \fill[white] (85.607,37.607) circle (.9);
  \fill[white] (75,12) circle (.75);
\end{scope}
